\documentclass[11pt]{article}

\usepackage[spanish, es-nodecimaldot]{babel}
\usepackage{amsmath}
\usepackage{amssymb}
\usepackage{graphicx}
\graphicspath{{./images/}}
\usepackage[utf8]{inputenc}
\usepackage[T1]{fontenc}
\usepackage[margin=1in]{geometry}
\usepackage{fancyhdr}
\usepackage{hyperref}

\hypersetup{
    pdftitle={Actividad 1.2 - Definición de la problemática},
    pdfauthor={Equipo 13 Proyectos: Felipe Colli, Agustin Concha, Joaquin Pérez, Benjamín Bascuñan, Diego Alvarado},
    colorlinks=true, 
    linkcolor=black,
    urlcolor=cyan
}

\pagestyle{fancy}
\fancyhf{}
\fancyfoot[C]{\footnotesize Equipo 13 Proyectos}
\fancyfoot[R]{\thepage}
\renewcommand{\headrulewidth}{0pt}

\title{Tarea 1\\Definición de la problemática}
\author{Equipo 13 Proyectos: \\ Felipe Colli, Agustin Concha,\\ Joaquin Pérez, Benjamín Bascuñan,\\ Diego Alvarado}
\date{\today}

\begin{document}

\maketitle
\thispagestyle{fancy}

\section{Problemáticas de interés}

Durante la conversación grupal se identificaron distintas problemáticas ambientales y sociales relacionadas con el desafío general del curso. El equipo consideró especialmente relevante analizar situaciones que afecten de forma directa a las personas y al entorno, y que además puedan ser abordadas mediante herramientas de ingeniería, diseño y tecnología.

Entre las problemáticas discutidas se encuentran las siguientes:

\begin{itemize}
    \item Acumulación de residuos en espacios públicos y naturales.
    \item Contaminación de playas y zonas costeras por basura de pequeño tamaño.
    \item Presencia de colillas de cigarro, plásticos y fragmentos de vidrio en la arena.
    \item Dificultades de los sistemas de limpieza tradicionales para retirar residuos mezclados o enterrados.
    \item Riesgos para visitantes, fauna y ecosistemas debido a residuos persistentes en el entorno.
\end{itemize}

Luego de analizar estas alternativas, el equipo decidió profundizar en la contaminación por residuos pequeños en playas, debido a su impacto ambiental, su relación con la seguridad de las personas y la dificultad de detectar este tipo de basura mediante métodos de limpieza convencionales.

\section{Investigación rápida}

Para comprender mejor la problemática seleccionada, se realizó una revisión inicial de antecedentes sobre contaminación costera y limpieza de playas. Esta búsqueda permitió identificar que los residuos presentes en la arena no corresponden únicamente a objetos grandes y visibles, sino también a elementos de menor tamaño, tales como colillas de cigarro, tapas, envoltorios, fragmentos de plástico y vidrio.

Una característica relevante de este problema es que parte de los residuos puede quedar cubierta por la arena debido al tránsito de personas, el viento y el movimiento de las mareas. Como consecuencia, una playa puede aparentar estar limpia en su superficie, aunque todavía contenga basura enterrada o mezclada con la arena.

Además, los métodos de limpieza manual suelen concentrarse en residuos visibles y de mayor tamaño. Esto dificulta la remoción constante de basura pequeña, especialmente en playas con alta afluencia de visitantes o en sectores donde los residuos se acumulan con rapidez.

La investigación preliminar también permitió reconocer que la contaminación de playas afecta la experiencia recreativa de las personas, puede generar riesgos físicos por la presencia de fragmentos cortantes y tiene consecuencias negativas para el ecosistema costero, ya que algunos residuos pueden ser ingeridos por animales o permanecer en el ambiente durante largos períodos.

\section{Problemática definida}

La problemática que se estudiará corresponde a la acumulación de residuos pequeños, visibles y enterrados, en la arena de las playas chilenas.

En particular, se observa que los métodos de limpieza tradicionales no siempre logran retirar de manera efectiva residuos como colillas de cigarro, fragmentos de plástico, tapas y vidrios que se encuentran mezclados con la arena. Esta situación puede afectar la seguridad de los visitantes, deteriorar la calidad del espacio recreativo y contribuir a la contaminación del ecosistema costero.

La problemática se formula de la siguiente manera:

\begin{quote}
En playas chilenas con alta afluencia de visitantes, las personas y las comunidades costeras enfrentan la persistencia de residuos pequeños en la arena, tanto en la superficie como enterrados, debido a las limitaciones de los métodos habituales de limpieza. Esto afecta la seguridad, la experiencia de uso de las playas y la conservación del entorno costero.
\end{quote}

\end{document}
